% !TEX root = ../main.tex

\chapter{实验平台与测量方法}

\section{轴向分级燃烧实验平台}

\subsection{系统总体结构搭建}

本实验所用轴向分级燃烧器由课题组自主设计。燃烧器由主燃室（一级燃烧区）、流动发展段和二级燃烧室（测试段）三部分组成，各部分以法兰连接，法兰连接面用石墨垫片进行密封。整个燃烧器安装在铝型材搭建的支撑框架上，二级燃烧室四面安装有耐高温石英玻璃观察窗。燃烧器出口连接有排烟段，将燃烧产物导至燃烧器外部。

\begin{figure}[H]
	\centering
	\includegraphics[width=0.9\textwidth]{fig1}
	\caption{轴向分级燃烧器整体结构示意图}
	\label{fig:burner_structure}
\end{figure}

气路供给系统包括两路独立的供气线路：空气路由螺杆压缩机提供，经过稳压罐、冷干机、PID闭环控制系统后送入燃烧器；燃料路由高压气瓶经减压阀、质量流量计后送入燃烧器。两路气体在进入燃烧器前相互独立，可以实现对空气和燃料流量的独立精确控制。

\subsection{主燃室结构}

主燃室采用多孔板燃烧方式。燃料通过与进气管相连接的细管注入到进气管中，大流量空气通过细管间的间隙腔体到达多孔板表面，两者在多孔板与钝体板之间的凹陷空间内完成混合，随后由下游值班火焰点燃。燃料细管出口与空气出口采用相互间隔的分布式布置方式，有利于燃料与空气在到达多孔板之前就进行预混合。

进气腔的孔隙结构数目为400个，直径为5mm。为促进一级燃料与空气的充分混合，设计加工了黄铜整流多孔板，孔径为1.2mm，孔数目为2607个，兼具均匀流速分布和防止火焰回火到进气腔的功能。

\begin{figure}[H]
	\centering
	\includegraphics[width=0.7\textwidth]{fig12}
	\caption{主燃室结构示意图}
	\label{fig:primary_combustor}
\end{figure}

\subsection{流动发展段与二级燃烧室}

流动发展段为渐缩矩形截面腔体，其主要功能是整流和加速：高温燃气经过渐缩段后流线趋于平行，速度分布更加均匀；截面面积减小使流速增加，符合轴向分级燃烧中二级燃烧区高速横流的工程特征。

\begin{figure}[H]
	\centering
	\includegraphics[width=0.8\textwidth]{fig13}
	\caption{流动发展段与二级燃烧室结构示意图}
	\label{fig:secondary_combustor}
\end{figure}

二级燃烧室截面尺寸为 \(115\mathrm{mm}\times 40\mathrm{mm}\) ，长度为 \(225\mathrm{mm}\) ，四面安装耐高温石英玻璃窗口，窗口内壁面与腔体内壁面平齐安装，以保证测试段内壁光滑、无阶梯状突起。二级燃烧室的上壁面安装有射流喷嘴，二级燃料经喷嘴垂直射入主流中，形成 \(90^{\circ}\) 横向射流。

\subsection{气路与控制系统}

空气路由室外螺杆压缩机供气，经稳压罐、冷干机后进入PID闭环控制回路。在调试过程中，发现上游压缩机持续补压导致管路压力波动，直接影响流量稳定性。针对这一问题，采用PID控制器实时采集质量流量计反馈信号，驱动快速响应比例调节阀自动调节开度，经多轮参数调整后，将流量波动控制在设定值 \(\pm 1.5\%\) 以内。燃料路由高压气瓶经两级减压阀（减压至0.3-0.5MPa）、质量流量计后送入燃烧器。点火系统采用高压电火花点火器，配备远程遥控开关，实现远程安全点火。

\subsection{数据采集系统开发}

所有流量及温度传感器通过RS-485通讯协议与计算机连接，采用LabVIEW平台开发了数据采集程序。程序中通过条件判断结构自动创建数据文件，发送十六位功能码实现瞬时流量和温度的精准读取，并设置了可调读数间隔。经过多次调试，解决了通讯延迟与数据丢帧问题，实现了多通道数据的同步采集、实时显示和自动存储，为后续实验奠定了技术基础。

\begin{figure}[H]
	\centering
	\includegraphics[width=0.8\textwidth]{fig2}
	\caption{LabVIEW数据采集程序界面}
	\label{fig:labview_interface}
\end{figure}

\section{收缩体的设计加工与安装}

收缩体是本课题实现“下游流道几何变化”的核心部件，用于改变流场结构。

\subsection{型线设计与计算}

收缩段型线设计需要注意型线的斜率应是渐变的，应保证气流均匀、平直，保证各部分平稳衔接等。目前关于收缩段设计尚无严格的空气动力学理论提供相应的指导，因此没有最优的理论设计型线。常用的收缩段设计方法均为经验公式。本文采用的为五次型线和Witoszynski曲线。五次曲线公式：

\begin{equation}
	\frac{D_i - D_t}{D - D_t} = 1 - 10\left(\frac{x}{L}\right)^3 +15\left(\frac{x}{L}\right)^4 -6\left(\frac{x}{L}\right)^5
\end{equation}

其中 \(D\) 为收缩段入口直径， \(D_{i}\) 为横坐标为 \(x\) 处的截面直径， \(D_{t}\) 为收缩段出口直径， \(L\) 为收缩距离。

Witoszynski曲线来源于风洞收缩段设计理论，使沿程速度梯度分布均匀，其表达式为：

\begin{equation}
	y(x) = \frac{y_0}{\sqrt{1 - \left[1 - \left(\frac{y_0}{y_i}\right)^2\right] \frac{(1 - x^2/L^2)^2}{(1 + x^2/3L^2)^3}}}
\end{equation}

其中 \(x\) 为沿轴向的距离， \(y(x)\) 为横坐标 \(x\) 处的截面半径， \(L\) 为收缩段长度， \(y_i\) 为入口半径， \(y_0\) 为出口半径。使用Python对用到的型线进行绘制。

\begin{figure}[H]
	\centering
	\includegraphics[width=0.7\textwidth]{fig3}
	\caption{收缩体型线对比}
	\label{fig:contour_profiles}
\end{figure}

\subsection{设计与加工}

将绘制的型线导入SolidWorks建立三维模型。对模型进行机加工，材料选用304不锈钢，加工精度控制在 \(\pm 0.02\mathrm{mm}\) 。加工完成后对表面进行打磨处理，按照单一变量加工了三个收缩体。

\begin{figure}[H]
	\centering
	\includegraphics[width=0.8\textwidth]{fig4}
	\caption{收缩体三维模型}
	\label{fig:contour_3d}
\end{figure}

\section{测量系统}

\subsection{温度测量}

温度测量使用B型热电偶（量程 \(0-1800^{\circ}C\) ，精度优于 \(\pm 0.5\%\) ），用于测量二级燃烧室入口及内部的高温气流温度，热电偶布置在二级燃烧室入口处，通过串口采集温度数据。

\subsection{流量测量}

所有气体流量使用热式质量流量计测量，通过针阀和球阀进行手动调节控制。流量计量程包括：一级空气2000 SLM、一级甲烷300 SLM、二级燃料20 SLM、二级空气50SLM，测量精度 \(\pm 1.0-1.5\%\) FS。对于甲烷和乙烯，需根据气体转换系数将空气标定读数转换为实际气体流量。

\subsection{紫外火焰探测器}

使用YJH-1E紫外火焰探测器（4-20mA输出）监测火焰信号。该探测器对185-260nm紫外辐射敏感，具有日盲特性，输出电流与火焰强度相关。但使用采集模块配套软件采集时，发现数据刷新速度慢，存储频率低，无法记录火焰的快速变化。

\begin{figure}[H]
	\centering
	\includegraphics[width=0.8\textwidth]{fig18}
	\caption{探测器敏感区域}
	\label{fig:detector_sensitivity}
\end{figure}

使用采集卡测量探测器的4-20mA输出端，发现探测器本身对火焰变化的响应速度远快于配套软件的采集存储频率——瓶颈在于采集模块的软件端。于是，改用Python编写串口读取程序，直接通过RS-485接口读取探测器数据。优化后，有效数据存储频率从原系统的5Hz提升至约20Hz，实现了4倍的采集存储速度。20Hz采样虽不足以进行波形细节分析，但已足以反映吹熄前信号的起伏趋势——数据中能明显看到信号在“有火”与“无火”阈值之间的交替变化，证实了吹熄前确实存在信号波动过程，目前可供分析的实验吹熄数据均来自于这一探测器。

然而，该Python高频读取方案受限于电脑硬件性能，且20Hz采样率仍无法获取高频的吹熄细节。因此，后续有必要从硬件层面探索新的采集方案。

\section{本章小结}

本章完成了以下具体工作：

（1）完成了轴向分级燃烧器主体结构的安装与气密性检查，使实验具备了基本运行条件。

（2）完成了收缩体的设计计算、SolidWorks建模和加工，完成了从Python型线绘制到实物加工的全过程。

（3）完成了气路系统PID闭环控制的调试，解决了压缩机补压引起的压力波动问题，将流量波动控制在设定值 \(\pm 1.5\%\) 以内。

（4）基于LabVIEW平台开发了多通道RS485数据采集程序，同时搭建了B型热电偶温度测试系统，完成了流量、温度的同步读取存储。

（5）完成了紫外火焰探测器的搭建与信号采集优化，通过Python高频串口读取将有效采样频率从5Hz提升至20Hz，证实了吹熄前低频信号波动过程。虽仍有局限，但已具备初步分析能力，目前实验吹熄数据均基于此探测器获取。

以上工作使实验具备了稳定运行、同步数据采集和火焰信号量化分析的基础条件。